% Oversells — nazorat ishi shabloni (XeLaTeX, exam klassi).
% Exam paper template based on the exam class (compile with XeLaTeX).
\documentclass[12pt,a4paper,addpoints]{exam}

\usepackage[margin=2cm]{geometry}
\usepackage{fontspec}
\setmainfont{Noto Serif}
\setsansfont{Noto Sans}
\usepackage{babel}
\babelprovide[import=uz-Latn, main]{uzbek}
\setlength{\emergencystretch}{3em}
\usepackage{amsmath, amssymb}

% exam klassi yozuvlarini oʻzbekchaga oʻgiramiz / Uzbek labels for the exam class
\pointpoints{ball}{ball}
\bonuspointpoints{qoʻshimcha ball}{qoʻshimcha ball}
\renewcommand{\solutiontitle}{\noindent\textbf{Yechim:}\enspace}
\hqword{Savol}
\hpword{Ball}
\hsword{Olingan}
\htword{Jami}
\vqword{Savol}
\vpword{Ball}
\vsword{Olingan}
\vtword{Jami}

\pagestyle{headandfoot}
\firstpageheader{Algebra, 8-sinf}{Nazorat ishi №\,2}{I variant}
\runningheader{Algebra, 8-sinf}{}{I variant}
\firstpagefooter{}{\thepage-sahifa}{}
\runningfooter{}{\thepage-sahifa}{}

% Yechimlarni koʻrsatish uchun quyidagi qatorni yoqing / uncomment to print solutions
% \printanswers

\begin{document}

\begin{center}
  {\Large\bfseries Kvadrat tenglamalar}\\[0.4em]
  Vaqt: 45 daqiqa \quad Jami: \numpoints\ ball
\end{center}

\vspace{0.5em}
\noindent Familiya, ism: \hrulefill\quad Sinf: \rule{3cm}{0.4pt}

\begin{center}
  \gradetable[h][questions]
\end{center}

\begin{questions}

\question[2] $x^2 - 9 = 0$ tenglamaning ildizlarini toping.
\begin{choices}
  \choice $x = 3$
  \choice $x = -3$
  \CorrectChoice $x = \pm 3$
  \choice ildizi yoʻq
\end{choices}

\question[2] $x^2 - 5x + 6 = 0$ tenglamaning diskriminanti nechaga teng?
\begin{oneparchoices}
  \choice $-1$
  \CorrectChoice $1$
  \choice $25$
  \choice $49$
\end{oneparchoices}

\question[3] Tenglamani yeching: $2x^2 + 3x - 2 = 0$.
\begin{solution}
  $D = 9 + 16 = 25$, demak $x_1 = \frac{-3 + 5}{4} = \frac12$, $x_2 = \frac{-3 - 5}{4} = -2$.
\end{solution}
\vspace{3cm}

\question[3] Toʻgʻri toʻrtburchakning yuzi $48\ \text{sm}^2$, perimetri esa $28$~sm.
Uning tomonlarini toping.
\vspace{3cm}

\end{questions}

\end{document}
